% Physical unit 024; accepted Session 02.
\GordanUnitStart{024}{printed page 19}
\GordanTopFolio{19}
is true, $U$ will be either positive or negative according as
\[
 {u\over \sin\varphi}\left\{{u\over \sin\varphi}-{u_2\over \sin\varphi_2}\right\}
\]
is either positive or negative.

Correspondingly, when
\[
 {u\over \sin\varphi}
\]
does not vanish again, $U$ ought to be positive; but when it has advanced beyond $0$, $U$ may be either positive or negative.

\begin{center}\S\ 15.\end{center}

Hence we state the theorem:
\begin{quote}
``When some point lying very near a geodetic line is joined by geodetic lines with the points $A$ and $B$, the sum of the lengths of those lines, if $\psi<\psi'$ is the case, is greater than the length `$A\ B$'; if $\psi>\psi$ is the case, it is either smaller or greater.''
\end{quote}

\bigskip\hrule\bigskip

\begin{center}
{\large CHAPTER V.}\\[0.8em]
\S\ 1.
\end{center}

Now, however, the shortest of all lines, whose definition we proposed at the beginning of the preceding chapter, must be discussed. For this purpose we shall consider all geodetic lines containing the point $A(l,\omega)$. Their equations are all obtained in turn by allowing in
\[
 \psi-\omega=\pm\int_{\varphi_1}^{\varphi}d\varphi\cos l_0{\sqrt{1-e^2+e^2\sin^2l_0\cos^2\varphi}\over 1-\sin^2l_0\cos^2\varphi}
\]
(in which the angles $\varphi$ and $\varphi_1$ are known to be defined by the formulae
\[
 \sin l_0\cos\varphi=\sin l\qquad \sin l_0\cos\varphi_1=\sin l_1
\]
) the constant $l_0$ to take all numbers from $0$ to ${\pi\over2}$; when $l_0$ has been fixed, however, the equation still remains double, since the integral can be taken as either positive or negative, so that two geodetic lines touch the same circle.

\begin{center}\S\ 2.\end{center}

There are four genera of these lines. For, when the point $\Pi$ is carried on any one of them to a neighboring point $(l,\psi)$, one has:
